

\section{Recursive TKF Models}

Another way to nest TKF models is to borrow from grammar theory and allow recursion.
This was previously used to develop RNA evolutionary models \cite{Holmes2004}, and is developed here in a general way for proteins
and genomic DNA.

We illustrate how the TKF-mixed domain model can be interpreted as a stochastic grammar,
 developing four examples---recursive protein domains (i.e. arbitrary nesting of motifs), a basic model of RNA foldback structure, a second more sophisticated model of RNA structure, and a basic model of a genome---that
 highlight the structure of such models as a series of stepwise grammar elaborations that constitute tree-adjoining moves on the space of grammars.

\subsection{Example One: Left-Recursive TKF (L-TKF)}

\newcommand\tok{s}
\newcommand\ntlinks{{\tt L}}
\newcommand\ntimm{{\tt K}}
\newcommand\ntmor{{\tt M}}
\newcommand\ntnew{{\tt N}}
\newcommand\ntsur{{\tt S}}
\newcommand\ntexp{{\tt E}}
\newcommand\ntaln{{\tt A}}
\newcommand\ntins{{\tt I}}
\newcommand\ntdel{{\tt D}}
\newcommand\ntparent{{\tt P}}
\newcommand\ntchild{{\tt C}}

\newcommand\ntfrag{{\tt F}}
\newcommand\ntgen{{\tt G}}
\newcommand\ntrem{{\tt R}}

\newcommand\ntbegin{{\tt B}}

\newcommand\nullability{\eta}

\newcommand\term{{\tt T}}


We can imagine a links model where, in a domain of type $\srcdom$,
each mortal link is either (with probability $\tok_\srcdom$) associated with a character,
or (with probability $(1 - \tok_\srcdom) \domdist_{\srcdom\destdom}$ associated with its own
independently-evolving links model of domain type $\destdom$.
Links models can thus be nested {\em ad infinitum}
\begin{eqnarray*}
\model_\srcdom & = &
\tkflinks(\mixture_{\sumidx\sim\tok}(L^{(\sumidx)}_\srcdom);\insrate_\srcdom,\delrate_\srcdom)
\\
L^{(1)}_\srcdom & = & \hmmproc(\{\exch^{(\class)},\eqm^{(\class)}\}_\class;\ext^{(\srcdom)},\classdist_\srcdom) \\
L^{(0)}_\srcdom & = & \mixture_{\destdom\sim\domdist_\destdom}(\model_\destdom)
\end{eqnarray*}
where $\hmmproc$ denotes the Markovian fragment process (Section~\ref{sec:mixdom}),
and $\mixture_{\sumidx\sim\tok}$ is defined for the Bernoulli index variable $\sumidx$ and probability $\tok$ as it was for categorical index variables
\[
\state \sim \mixture_{\sumidx\sim p}(\model(\theta_\sumidx);p)\ \Leftrightarrow\ 
\sumidx \sim \berndist(p),\ \state \sim \model(\theta_\sumidx)
\]

We postpone the TKF92-like augmentations of $L^{(1)}_\srcdom$ (fragment types $\frag$ and site classes $\class$) for now, and start with a simplified TKF91-like version
that allows full recursively-nested domains but allows only single-character fragments with one site class per domain
\begin{eqnarray*}
\model'_\srcdom & = &
\tkflinks(\mixture_{\sumidx\sim\tok}(L'^{(\sumidx)}_\srcdom);\insrate_\srcdom,\delrate_\srcdom)
\\
L'^{(1)}_\srcdom & = & \subproc(\exch_\srcdom,\eqm_\srcdom) \\
L'^{(0)}_\srcdom & = & \mixture_{\destdom\sim\domdist_\srcdom}(\model'_\destdom)
\end{eqnarray*}
As with the nested TKF HMM, we have to account for the probability of zero-length components, leading to null cycles during likelihood and inference computations.

The joint distribution over ancestor-descendant alignments under the recursive TKF model
is described by the following stochastic context-free grammar (SCFG)
\[
\begin{array}{lrcccl}
\mbox{Symbol interpretation} & \mbox{LHS} & & \multicolumn{2}{c}{\mbox{RHS}} & \mbox{Probability} \\
\hline
\mbox{\underline{\bf L}ink sequence, domain $\srcdom$} &
\ntlinks_\srcdom & \to & \ntlinks_\srcdom & \ntmor_\srcdom & \kappa_\srcdom \\
& & | & \ntimm_\srcdom & & 1 - \kappa_\srcdom \\
\mbox{Immortal lin\underline{\bf k}} &
\ntimm_\srcdom & \to & \ntnew_\srcdom & & \beta_\srcdom \\
& & | & \epsilon & & 1 - \beta_\srcdom \\
\mbox{\underline{\bf M}ortal link} &
\ntmor_\srcdom & \to & \ntsur_\srcdom & & \alpha_\srcdom \\
& & | & \ntexp_\srcdom & & 1 - \alpha_\srcdom \\
\mbox{\underline{\bf S}urviving mortal link} &
\ntsur_\srcdom & \to & \ntaln_\srcdom & \ntnew_\srcdom & \beta_\srcdom \\
& & | & \ntaln_\srcdom & & 1 - \beta_\srcdom \\
\mbox{\underline{\bf E}xpired mortal link} &
\ntexp_\srcdom & \to & \ntdel_\srcdom & \ntnew_\srcdom & \gamma_\srcdom \\
& & | & \ntdel_\srcdom & & 1 - \gamma_\srcdom \\
\mbox{\underline{\bf N}ewborn mortal link(s)} &
\ntnew_\srcdom & \to & \ntins_\srcdom & \ntnew_\srcdom & \beta_\srcdom \\
& & | & \ntins_\srcdom & & 1 - \beta_\srcdom \\
\mbox{\underline{\bf A}ligned component} &
\ntaln_\srcdom & \to & \term_{\anctok\destok} & & \tok_\srcdom \eqm_{\srcdom\anctok} \exp(\exch_\srcdom \evoltime)_{\anctok\destok}\\
& & | & \ntlinks_\destdom & & (1 - \tok_\srcdom) \domdist_{\srcdom\destdom} \\
\mbox{\underline{\bf I}nserted component} &
\ntins_\srcdom & \to & \term_{\gap\destok} & & \tok_\srcdom \eqm_{\srcdom\destok} \\
& & | & \ntchild_\destdom & & (1 - \tok_\srcdom) \domdist_{\srcdom\destdom} \\
\mbox{\underline{\bf D}eleted component} &
\ntdel_\srcdom & \to & \term_{\anctok\gap} & & \tok_\srcdom \eqm_{\srcdom\anctok} \\
& & | & \ntparent_\destdom & & (1 - \tok_\srcdom) \domdist_{\srcdom\destdom} \\
\mbox{\underline{\bf C}hild link sequence (inserted)} &
\ntchild_\srcdom & \to & \ntchild_\srcdom & \ntins_\srcdom & \kappa_\srcdom \\
& & | & \epsilon & & 1 - \kappa_\srcdom \\
\mbox{\underline{\bf P}arent link sequence (deleted)} &
\ntparent_\srcdom & \to & \ntparent_\srcdom & \ntdel_\srcdom & \kappa_\srcdom  \\
& & | & \epsilon & & 1 - \kappa_\srcdom \\
\end{array}
\]

The standard path from here is to transform the grammar to Chomsky Normal Form  \cite{stolcke-1995-efficient}.
We don't need to go all the way down that path;
we only need to remove $\epsilon$-productions and null cycles.

In order to train by EM, every time we remove $\epsilon$-productions and null cycles,
we need to be able to convert a rule count in the $\epsilon$-eliminated grammar
to a set of rule counts in the original grammar.
For each non-nullable nonterminal $\xstate'$, when a rule
$\xstate' \to \ystate'$ appears (arising from an original bifurcation
$\xstate \to \ystate\;\zstate$ where $\zstate$ was nullable),
the expected count $c(\xstate' \to \ystate')$ should contribute
$c(\xstate' \to \ystate') \cdot \nullability(\zstate)$ to the count
of the original rule $\xstate \to \ystate\;\zstate$, and symmetrically
when the first child was nullable.

More precisely, for the nullability fixed-point iteration, this means
tracking an additional set of \emph{expected counts} alongside
the nullabilities.
For each nonterminal $\xstate$ with nullability $\nullability(\xstate)$,
define $\bar{c}(\xstate \to \alpha)$ as the expected number of times
rule $\xstate \to \alpha$ would have been used in the original grammar,
conditioned on $\xstate$ generating the empty string.
These satisfy analogous fixed-point equations:
\[
\bar{c}(\xstate \to \ystate_1 \cdots \ystate_k) =
\WF(\xstate \to \ystate_1 \cdots \ystate_k) \prod_{i=1}^k \nullability(\ystate_i)
+ \sum_{i=1}^k \frac{\WF(\xstate \to \ystate_1 \cdots \ystate_k) \prod_{j=1}^k \nullability(\ystate_j)}
{\nullability(\xstate)} \sum_{\alpha'} \bar{c}(\ystate_i \to \alpha')
\]
and can be iterated to convergence alongside the nullabilities.
Given posterior counts from the Inside-Outside algorithm on the
$\epsilon$-eliminated grammar, the original-grammar rule counts
are recovered by:
\begin{enumerate}
\item For each non-null production $\xstate' \to \alpha'$ in the
  $\epsilon$-eliminated grammar, its posterior count $c'(\xstate' \to \alpha')$
  contributes directly to the corresponding original rule.
\item For each ``nullability shortcut'' production
  $\xstate' \to \ystate'$ (arising from
  $\xstate \to \ystate\;\zstate$ with $\zstate$ nullable),
  the count $c'(\xstate' \to \ystate')$ contributes
  $c'(\xstate' \to \ystate')$ to the original rule
  $\xstate \to \ystate\;\zstate$, plus
  $c'(\xstate' \to \ystate') \cdot \bar{c}(\zstate \to \alpha)$
  to each rule within $\zstate$'s null derivation subtree.
\end{enumerate}
The general theory of $\epsilon$-elimination with EM count recovery,
including the null closure $(I - T_{\mathcal{Z}\mathcal{Z}})^{-1}$
and null cycle removal, is developed in
Appendix~\ref{sec:appendix-grammar} (Section~\ref{sec:null-management-appendix})

First we find the nullability
$\nullability(\xstate)$ of each nonterminal $\xstate$.
The nullability is the probability that a parse tree rooted in that nonterminal
yields the empty string.
We can't solve for these in closed form (at least not in the general recursive
model, where a link sequence can contain another link sequence of the same type).
Instead we can solve approximately by iterating towards a fixed point.

We first observe that the nullabilities collectively satisfy the following
\begin{eqnarray*}
    \nullability(\ntlinks_\srcdom) & = & \nullability(\ntimm_\srcdom) \frac{1 - \kappa_\srcdom}{1 - \kappa_\srcdom \nullability(\ntmor_\srcdom)} \\
    \nullability(\ntimm_\srcdom) & = & 1 - \beta_\srcdom + \beta_\srcdom \nullability(\ntnew_\srcdom) \\
    \nullability(\ntmor_\srcdom) & = & \alpha_\srcdom \nullability(\ntsur_\srcdom) + (1 - \alpha_\srcdom) \nullability(\ntexp_\srcdom) \\
    \nullability(\ntsur_\srcdom) & = & \nullability(\ntaln_\srcdom) \left( \beta_\srcdom \nullability(\ntnew_\srcdom) + (1 - \beta_\srcdom) \right) \\
    \nullability(\ntexp_\srcdom) & = & \nullability(\ntdel_\srcdom) \left( \gamma_\srcdom \nullability(\ntnew_\srcdom) + (1 - \gamma_\srcdom) \right) \\
    \nullability(\ntnew_\srcdom) & = & \nullability(\ntins_\srcdom) \frac{1 - \beta_\srcdom}{1 - \beta_\srcdom \nullability(\ntins_\srcdom)}  \\
    \nullability(\ntaln_\srcdom) & = & (1 - \tok_\srcdom) \sum_\destdom \domdist_{\srcdom\destdom} \nullability(\ntlinks_\destdom) \\
    \nullability(\ntins_\srcdom) & = & (1 - \tok_\srcdom)\sum_\destdom \domdist_{\srcdom\destdom} \nullability(\ntchild_\destdom) \\
    \nullability(\ntdel_\srcdom) & = & \nullability(\ntins_\srcdom) \\
    \nullability(\ntchild_\srcdom) & = & \frac{1 - \kappa_\srcdom}{1 - \kappa_\srcdom \nullability(\ntins_\srcdom)} \\
    \nullability(\ntparent_\srcdom) & = & \nullability(\ntchild_\srcdom) \\    
\end{eqnarray*}

A general procedure is (i) iterate to solve for the $\nullability(\ntchild_\srcdom)$ (substituting in the downstream definition of $\nullability(\ntins_\srcdom)$ so the formula becomes self-referential);
(ii) this directly yields $\nullability(\ntparent_\srcdom), \nullability(\ntins_\srcdom), \nullability(\ntdel_\srcdom), \nullability(\ntnew_\srcdom), \nullability(\ntexp_\srcdom), \nullability(\ntimm_\srcdom)$;
(iii) iterate to solve for $\nullability(\ntlinks_\srcdom)$ (again, first substituting to make it self-referential);
(iv) this directly yields the remaining $\nullability(\ntmor_\srcdom), \nullability(\ntsur_\srcdom), \nullability(\ntaln_\srcdom)$.

In detail: initialize $x_\srcdom^{(0)} \leftarrow 0$ for $\srcdom \in \ndom$.
Iterate to convergence
\[
    x_\srcdom^{(\sumidx+1)}
    \leftarrow
    \frac{1 - \kappa_\srcdom}{1 - \kappa_\srcdom (1 - \tok_\srcdom)\sum_\destdom \domdist_{\srcdom\destdom} x_\destdom^{(\sumidx)}}
\]
We then set $\nullability(\ntchild_\srcdom) \leftarrow \lim_{\sumidx \to \infty} x_\srcdom^{(\sumidx)}$
and set $\nullability(\ntparent_\srcdom), \nullability(\ntins_\srcdom), \nullability(\ntdel_\srcdom), \nullability(\ntnew_\srcdom), \nullability(\ntexp_\srcdom), \nullability(\ntimm_\srcdom)$ using the above equations.
Now set $y_\srcdom^{(0)} \leftarrow 0$ for $\srcdom \in \ndom$
and again iterate
\[
    y_\srcdom^{(\sumidx+1)} \leftarrow \frac{(1 - \kappa_\srcdom) \nullability(\ntimm_\srcdom)}{1 - \kappa_\srcdom (1 - \alpha_\srcdom) \nullability(\ntexp_\srcdom) 
      - \kappa_\srcdom \alpha_\srcdom  
      \left( 1 - \beta_\srcdom(1 - \nullability(\ntnew_\srcdom)) \right)
     (1 - \tok_\srcdom) \sum_\destdom \domdist_{\srcdom\destdom} y_\destdom^{(\sumidx)}
      }
\]
Then set
$\nullability(\ntlinks_\srcdom) \leftarrow \lim_{\sumidx \to \infty} y_\srcdom^{(\sumidx)}$
and set the remaining $\nullability(\ntmor_\srcdom), \nullability(\ntsur_\srcdom), \nullability(\ntaln_\srcdom)$ using the above equations.

We next develop a ``non-nullable'' version of the grammar that yields the same Inside
probabilities, but explicitly separates out $\epsilon$-generations.
For every nonterminal $\xstate_\srcdom$ we create a new nonterminal $\xstate'_\srcdom$
with rules that (by construction) never generate empty parse trees,
but are otherwise identical to those transforming $\xstate_\srcdom$.
In cases where the original grammars has bifurcation rules $\xstate_i \to \xstate_j \xstate_k$,
we need to introduce transitions $\xstate'_i \to \xstate'_j$ and $\xstate'_i \to \xstate'_k$,
to account for the missing nullability of $\xstate_j$ and $\xstate_k$.
(We can also eliminate $\ntimm_\srcdom$, which only has one outgoing rule
when its $\epsilon$-production is removed, and fold $\ntlinks_\srcdom\to\ntimm_\srcdom\to\ntmor_\srcdom$ into $\ntlinks_\srcdom\to\ntmor_\srcdom$.)

\[
\begin{array}{lrcccl}
\mbox{Symbol interpretation} & \mbox{LHS} & & \multicolumn{2}{c}{\mbox{RHS}} & \mbox{Probability} \\
\hline
\mbox{\underline{\bf L}ink sequence, domain $\srcdom$} &
\ntlinks'_\srcdom & \to & \ntlinks'_\srcdom & \ntmor'_\srcdom & \kappa_\srcdom \\
\mbox{(nonempty)} & & | & \ntlinks'_\srcdom & & \kappa_\srcdom \nullability(\ntmor_\srcdom) \\
& & | & \ntmor'_\srcdom & & \kappa_\srcdom \nullability(\ntlinks_\srcdom) \\
& & | & \ntnew'_\srcdom & & (1 - \kappa_\srcdom) \beta_\srcdom \\
\mbox{\underline{\bf M}ortal link (nonempty)} &
\ntmor'_\srcdom & \to & \ntsur'_\srcdom & & \alpha_\srcdom \\
& & | & \ntexp'_\srcdom & & 1 - \alpha_\srcdom \\
\mbox{\underline{\bf S}urviving mortal link} &
\ntsur'_\srcdom & \to & \ntaln'_\srcdom & \ntnew'_\srcdom & \beta_\srcdom \\
\mbox{(etc.; all $\xstate'_\srcdom$ are nonempty)} & & | & \ntnew'_\srcdom & & \beta_\srcdom \nullability(\ntaln_\srcdom) \\
& & | & \ntaln'_\srcdom & & 1 - \beta_\srcdom(1 - \nullability(\ntnew_\srcdom)) \\
\mbox{\underline{\bf E}xpired mortal link} &
\ntexp'_\srcdom & \to & \ntdel'_\srcdom & \ntnew'_\srcdom & \gamma_\srcdom \\
& & | & \ntnew'_\srcdom & & \gamma_\srcdom \nullability(\ntdel_\srcdom) \\
& & | & \ntdel'_\srcdom & & 1 - \gamma_\srcdom(1 - \nullability(\ntnew_\srcdom)) \\
\mbox{\underline{\bf N}ewborn mortal link(s)} &
\ntnew'_\srcdom & \to & \ntins'_\srcdom & \ntnew'_\srcdom & \beta_\srcdom \\
& & | & \ntnew'_\srcdom & & \beta_\srcdom \nullability(\ntins_\srcdom) \\
& & | & \ntins'_\srcdom & & 1 - \beta_\srcdom(1 - \nullability(\ntnew_\srcdom)) \\
\mbox{\underline{\bf A}ligned component} &
\ntaln'_\srcdom & \to & \term_{\anctok\destok} & & \tok_\srcdom \eqm_{\srcdom\anctok} \exp(\exch_\srcdom \evoltime)_{\anctok\destok}\\
& & | & \ntlinks'_\destdom & & (1 - \tok_\srcdom) \domdist_{\srcdom\destdom} \\
\mbox{\underline{\bf I}nserted component} &
\ntins'_\srcdom & \to & \term_{\gap\destok} & & \tok_\srcdom \eqm_{\srcdom\destok} \\
& & | & \ntchild'_\destdom & & (1 - \tok_\srcdom) \domdist_{\srcdom\destdom} \\
\mbox{\underline{\bf D}eleted component} &
\ntdel'_\srcdom & \to & \term_{\anctok\gap} & & \tok_\srcdom \eqm_{\srcdom\anctok} \\
& & | & \ntparent'_\destdom & & (1 - \tok_\srcdom) \domdist_{\srcdom\destdom} \\
\mbox{\underline{\bf C}hild link sequence (inserted)} &
\ntchild'_\srcdom & \to & \ntchild'_\srcdom & \ntins'_\srcdom & \kappa_\srcdom \\
& & | & \ntchild'_\srcdom & & \kappa_\srcdom \nullability(\ntins_\srcdom) \\
& & | & \ntins'_\srcdom & & \kappa_\srcdom \nullability(\ntchild_\srcdom) \\
\mbox{\underline{\bf P}arent link sequence (deleted)} &
\ntparent'_\srcdom & \to & \ntparent'_\srcdom & \ntdel'_\srcdom & \kappa_\srcdom  \\
& & | & \ntparent'_\srcdom & & \kappa_\srcdom \nullability(\ntdel_\srcdom) \\
& & | & \ntdel'_\srcdom & & \kappa_\srcdom \nullability(\ntparent_\srcdom) \\
\end{array}
\]

The final step is to remove null cycles: chains of unit productions resulting in the same nonterminal.
Specifically we need to remove $\ntlinks'_\srcdom\to\ntmor'_\srcdom\to\ntsur'_\srcdom\to\ntaln'_\srcdom\to\ntlinks'_\destdom$,
$\ntchild'_\srcdom\to\ntins'_\srcdom\to\ntchild'_\destdom$ and $\ntparent'_\srcdom\to\ntdel'_\srcdom\to\ntparent'_\destdom$.
We do this by deleting the $\ntaln'_\srcdom\to\ntlinks'_\destdom$, $\ntins'_\srcdom\to\ntchild'_\destdom$ and $\ntdel'_\srcdom\to\ntparent'_\destdom$
transitions, adding compensatory self-loops and self-bifurcations to $\ntlinks'_\srcdom$, $\ntparent'_\srcdom$, and $\ntchild'_\srcdom$
to account for the now-broken paths through $\ntaln'_\srcdom$, $\ntins'_\srcdom$ and $\ntdel'_\srcdom$.
The modified grammar is

\[
\begin{array}{lrcccl}
\mbox{Symbol} & \mbox{LHS} & & \multicolumn{2}{c}{\mbox{RHS}} & \mbox{Probability} \\
\hline
\mbox{\underline{\bf L}ink} &
\ntlinks''_\srcdom & \to & \ntlinks''_\srcdom & \ntmor''_\srcdom & \kappa_\srcdom \\
& & | & \ntlinks''_\srcdom & \ntlinks''_\destdom & \kappa_\srcdom \alpha_\srcdom \left( 1 - \beta_\srcdom(1 - \nullability(\ntnew_\srcdom)) \right) (1 - \tok_\srcdom) \domdist_{\srcdom\destdom} \\
& & | & \ntlinks''_\destdom & & \ltrans_{\srcdom\destdom} \\
& & | & \ntmor''_\srcdom & & \kappa_\srcdom \nullability(\ntlinks_\srcdom) \\
& & | & \ntnew''_\srcdom & & (1 - \kappa_\srcdom) \beta_\srcdom \\
\mbox{\underline{\bf M}ortal} &
\ntmor''_\srcdom & \to & \ntsur''_\srcdom & & \alpha_\srcdom \\
& & | & \ntexp''_\srcdom & & 1 - \alpha_\srcdom \\
\mbox{\underline{\bf S}urviving} &
\ntsur''_\srcdom & \to & \ntaln''_\srcdom & \ntnew''_\srcdom & \beta_\srcdom \\
& & | & \ntnew''_\srcdom & & \beta_\srcdom \nullability(\ntaln_\srcdom) \\
& & | & \ntaln''_\srcdom & & 1 - \beta_\srcdom(1 - \nullability(\ntnew_\srcdom)) \\
\mbox{\underline{\bf E}xpired} & \ntexp''_\srcdom & \to & \ntdel''_\srcdom & \ntnew''_\srcdom & \gamma_\srcdom \\
& & | & \ntparent''_\destdom & \ntnew''_\srcdom & \gamma_\srcdom (1 - \tok_\srcdom) \domdist_{\srcdom\destdom} \\
& & | & \ntnew''_\srcdom & & \gamma_\srcdom \nullability(\ntdel_\srcdom) \\
& & | & \ntdel''_\srcdom & & 1 - \gamma_\srcdom(1 - \nullability(\ntnew_\srcdom)) \\
& & | & \ntparent''_\destdom & & \left( 1 - \gamma_\srcdom(1 - \nullability(\ntnew_\srcdom)) \right) (1 - \tok_\srcdom) \domdist_{\srcdom\destdom} \\
\mbox{\underline{\bf N}ewborns} & \ntnew''_\srcdom & \to & \ntins''_\srcdom & \ntnew''_\srcdom & \beta_\srcdom \\
& & | & \ntchild''_\destdom & \ntnew''_\srcdom & \beta_\srcdom (1 - \tok_\srcdom) \domdist_{\srcdom\destdom} \\
& & | & \ntnew''_\srcdom & & \beta_\srcdom \nullability(\ntins_\srcdom) \\
& & | & \ntins''_\srcdom & & 1 - \beta_\srcdom(1 - \nullability(\ntnew_\srcdom)) \\
& & | & \ntchild''_\destdom & & \left( 1 - \beta_\srcdom(1 - \nullability(\ntnew_\srcdom)) \right) (1 - \tok_\srcdom) \domdist_{\srcdom\destdom} \\
\mbox{\underline{\bf A}ligned} & \ntaln''_\srcdom & \to & \term_{\anctok\destok} & & \tok_\srcdom \eqm_{\srcdom\anctok} \exp(\exch_\srcdom \evoltime)_{\anctok\destok}\\
\mbox{\underline{\bf I}nserted} & \ntins''_\srcdom & \to & \term_{\gap\destok} & & \tok_\srcdom \eqm_{\srcdom\destok} \\
\mbox{\underline{\bf D}eleted} & \ntdel''_\srcdom & \to & \term_{\anctok\gap} & & \tok_\srcdom \eqm_{\srcdom\anctok} \\
\mbox{\underline{\bf C}hild} & \ntchild''_\srcdom & \to & \ntchild''_\srcdom & \ntins''_\srcdom & \kappa_\srcdom \\
& & | & \ntchild''_\srcdom & \ntchild''_\destdom & \kappa_\srcdom (1 - \tok_\srcdom) \domdist_{\srcdom\destdom} \\
& & | & \ntchild''_\destdom & & \ctrans_{\srcdom\destdom} \\
& & | & \ntins''_\srcdom & & \kappa_\srcdom \nullability(\ntchild_\srcdom) \\
\mbox{\underline{\bf P}arent} &
\ntparent''_\srcdom & \to & \ntparent''_\srcdom & \ntdel''_\srcdom & \kappa_\srcdom  \\
& & | & \ntparent''_\srcdom & \ntparent''_\destdom & \kappa_\srcdom (1 - \tok_\srcdom) \domdist_{\srcdom\destdom} \\
& & | & \ntparent''_\destdom & & \ctrans_{\srcdom\destdom} \\
& & | & \ntdel''_\srcdom & & \kappa_\srcdom \nullability(\ntchild_\srcdom) \\
\end{array}
\]
where $\ltrans,\ctrans$ are the $\ndom\times\ndom$ transition matrices
\begin{eqnarray*}
    \ltrans_{\srcdom\destdom} & = & \kappa_\srcdom \left(
    \nullability(\ntmor_\srcdom) \delta_{\srcdom\destdom}  % Lm->Lm
+ \nullability(\ntlinks_\srcdom) % Lm->Mm
  \alpha_\srcdom % Mm->Sm
  \left(1 - \beta_\srcdom\left(1 - \nullability(\ntnew_\srcdom)\right)\right)  % Sm->Am
  \left(1 - \tok_\srcdom\right) \domdist_{\srcdom\destdom} \right)
  \\
    \ctrans_{\srcdom\destdom} & = & \kappa_\srcdom \left(
  \nullability(\ntins_\srcdom) \delta_{\srcdom\destdom}  % Cm->Cm
  + \nullability(\ntchild_\srcdom) % Cm->Im
  (1 - \tok_\srcdom) \domdist_{\srcdom\destdom}  % Im->Cn
  \right)
\end{eqnarray*}

Letting $\ltransinv=(I-\ltrans)^{-1}$ and $\ctransinv=(I-\ctrans)^{-1}$ be the geometric series sums of these matrices,
update the rules to sum over all
$\ntlinks''_\srcdom\to\ntlinks''_\destdom$,
$\ntchild''_\srcdom\to\ntchild''_\destdom$, and
$\ntparent''_\srcdom\to\ntparent''_\destdom$ transition chains explicitly
\[
\begin{array}{lrcccl}
\mbox{Symbol} & \mbox{LHS} & & \multicolumn{2}{c}{\mbox{RHS}} & \mbox{Probability} \\
\hline
\mbox{\underline{\bf L}ink} &
\ntlinks'''_\srcdom & \to & \ntlinks'''_\destdom & \ntmor'''_\destdom & \ltransinv_{\srcdom\destdom} \kappa_\destdom \\
& & | & \ntlinks'''_\destdom & \ntlinks'''_\bifdom & \ltransinv_{\srcdom\destdom} \kappa_\destdom \alpha_\destdom \left( 1 - \beta_\destdom(1 - \nullability(\ntnew_\destdom)) \right) (1 - \tok_\destdom) \domdist_{\destdom\bifdom} \\
& & | & \ntmor'''_\destdom & & \ltransinv_{\srcdom\destdom} \kappa_\destdom \nullability(\ntlinks_\destdom) \\
& & | & \ntnew'''_\destdom & & \ltransinv_{\srcdom\destdom} (1 - \kappa_\destdom) \beta_\destdom \\
\mbox{\underline{\bf C}hild} & \ntchild'''_\srcdom & \to & \ntchild'''_\destdom & \ntins'''_\destdom & \ctransinv_{\srcdom\destdom} \kappa_\destdom \\
& & | & \ntchild'''_\destdom & \ntchild'''_\bifdom & \ctransinv_{\srcdom\destdom} (1 - \tok_\destdom) \domdist_{\destdom\bifdom} \\
& & | & \ntins'''_\destdom & & \ctransinv_{\srcdom\destdom} \kappa_\destdom \nullability(\ntchild_\destdom) \\
\mbox{\underline{\bf P}arent} & \ntparent'''_\srcdom & \to & \ntparent'''_\destdom & \ntdel'''_\srcdom & \ctransinv_{\srcdom\destdom} \kappa_\destdom \\
& & | & \ntparent'''_\destdom & \ntparent'''_\bifdom & \ctransinv_{\srcdom\destdom} (1 - \tok_\destdom) \domdist_{\destdom\bifdom} \\
& & | & \ntdel'''_\destdom & & \ctransinv_{\srcdom\destdom} \kappa_\destdom \nullability(\ntchild_\destdom) \\
\end{array}
\]



Rules for $\xstate'''_\srcdom\ldots$ where $\xstate \in \{ \ntmor, \ntsur, \ntexp, \ntnew \}$ are just copied over from the corresponding $\xstate''_\srcdom\to\ldots$ rules,
changing $\xstate''$ to $\xstate'''$ on the right-hand side as well.


\[
\begin{array}{lrcccl}
\mbox{Symbol} & \mbox{LHS} & & \multicolumn{2}{c}{\mbox{RHS}} & \mbox{Probability} \\
\hline
\mbox{\underline{\bf M}ortal} & \ntmor'''_\srcdom & \to & \ntsur'''_\srcdom & & \alpha_\srcdom \\
& & | & \ntexp'''_\srcdom & & 1 - \alpha_\srcdom \\
\mbox{\underline{\bf S}urviving} & \ntsur'''_\srcdom & \to & \ntaln'''_\srcdom & \ntnew'''_\srcdom & \beta_\srcdom \\
& & | & \ntnew'''_\srcdom & & \beta_\srcdom \nullability(\ntaln_\srcdom) \\
& & | & \ntaln'''_\srcdom & & 1 - \beta_\srcdom(1 - \nullability(\ntnew_\srcdom)) \\
\mbox{\underline{\bf E}xpired} & \ntexp'''_\srcdom & \to & \ntdel'''_\srcdom & \ntnew'''_\srcdom & \gamma_\srcdom \\
& & | & \ntparent'''_\destdom & \ntnew'''_\srcdom & \gamma_\srcdom (1 - \tok_\srcdom) \domdist_{\srcdom\destdom} \\
& & | & \ntnew'''_\srcdom & & \gamma_\srcdom \nullability(\ntdel_\srcdom) \\
& & | & \ntdel'''_\srcdom & & 1 - \gamma_\srcdom(1 - \nullability(\ntnew_\srcdom)) \\
& & | & \ntparent'''_\destdom & & \left( 1 - \gamma_\srcdom(1 - \nullability(\ntnew_\srcdom)) \right) (1 - \tok_\srcdom) \domdist_{\srcdom\destdom} \\
\mbox{\underline{\bf N}ewborns} & \ntnew'''_\srcdom & \to & \ntins'''_\srcdom & \ntnew'''_\srcdom & \beta_\srcdom \\
& & | & \ntchild'''_\destdom & \ntnew'''_\srcdom & \beta_\srcdom (1 - \tok_\srcdom) \domdist_{\srcdom\destdom} \\
& & | & \ntnew'''_\srcdom & & \beta_\srcdom \nullability(\ntins_\srcdom) \\
& & | & \ntins'''_\srcdom & & 1 - \beta_\srcdom(1 - \nullability(\ntnew_\srcdom)) \\
& & | & \ntchild'''_\destdom & & \left( 1 - \beta_\srcdom(1 - \nullability(\ntnew_\srcdom)) \right) (1 - \tok_\srcdom) \domdist_{\srcdom\destdom} \\
\end{array}
\]
We use this last opportunity to reintroduce mixtures of site and fragment classes

\[
\begin{array}{lrcccl}
\mbox{Symbol} & \mbox{LHS} & & \multicolumn{2}{c}{\mbox{RHS}} & \mbox{Probability} \\
\hline
\mbox{\underline{\bf A}ligned} & \ntaln'''_\srcdom & \to & \ntfrag_\frag & & \tok_\srcdom \fragdist_{\srcdom\frag} \\
\mbox{\underline{\bf I}nserted} & \ntins'''_\srcdom & \to & \ntgen_\frag & & \tok_\srcdom \fragdist_{\srcdom\frag} \\
\mbox{\underline{\bf D}eleted} & \ntdel'''_\srcdom & \to & \ntrem_\frag & & \tok_\srcdom \fragdist_{\srcdom\frag} \\
\mbox{\underline{\bf F}ragment} & \ntfrag_\frag & \to & \term_{\anctok\destok} & \ntfrag_\destfrag & \ext^{(\srcdom)}_{\frag\destfrag} \sum_\class \classdist_{\srcdom\frag\class} \eqm^{(\class)}_\anctok \exp(\revsub^{(\class)} \evoltime)_{\anctok\destok}\\
& & | & \term_{\anctok\destok} & & \notext^{(\srcdom)}_\frag \sum_\class \classdist_{\srcdom\frag\class} \eqm^{(\class)}_\anctok \exp(\revsub^{(\class)} \evoltime)_{\anctok\destok} \\
\mbox{\underline{\bf G}enerated} & \ntgen_\frag & \to & \term_{\gap\destok} & \ntgen_\destfrag & \ext^{(\srcdom)}_{\frag\destfrag} \sum_\class \classdist_{\srcdom\frag\class} \eqm^{(\class)}_\destok \\
& & | & \term_{\gap\destok} & & \notext^{(\srcdom)}_\frag \sum_\class \classdist_{\srcdom\frag\class} \eqm^{(\class)}_\destok \\
\mbox{\underline{\bf R}emoved} & \ntrem_\frag & \to & \term_{\anctok\gap} & \ntrem_\destfrag & \ext^{(\srcdom)}_{\frag\destfrag} \sum_\class \classdist_{\srcdom\frag\class} \eqm^{(\class)}_\anctok \\
& & | & \term_{\anctok\gap} & & \notext^{(\srcdom)}_\frag \sum_\class \classdist_{\srcdom\frag\class} \eqm^{(\class)}_\anctok \\
\end{array}
\]
and add the top-level rule:
\[
\begin{array}{lrcccl}
\mbox{Symbol} & \mbox{LHS} & & \multicolumn{2}{c}{\mbox{RHS}} & \mbox{Probability} \\
\hline
\mbox{\underline{\bf B}egin} & \ntbegin & \to & \ntlinks'''_1 & & 1 \\
& & | & \epsilon & & \nullability(\ntlinks_1) \\
\end{array}
\]

With no $\epsilon$-productions and a strictly acyclic topological transition sort order on nonterminals
of $\ntbegin \to \ntlinks \to \ntmor \to ( \ntsur \to (\ntaln,\ \ntnew \to \ntins \to \ntchild),\ \ntexp \to \ntdel \to \ntparent)$,
this grammar is ready for an Inside parser.

At this point our model is the desired
\begin{eqnarray*}
\model_\srcdom & = &
% Flagging this for review: why \sumidx\sim\tok ? \tok is not a distribution ?
\tkflinks(\mixture_{\sumidx\sim\tok}(L^{(\sumidx)}_\srcdom);\insrate_\srcdom,\delrate_\srcdom)
\\
L^{(1)}_\srcdom & = & \hmmproc(\{\exch^{(\class)},\eqm^{(\class)}\}_\class;\ext^{(\srcdom)},\classdist_\srcdom) \\
L^{(0)}_\srcdom & = & \mixture_{\destdom\sim\domdist_\destdom}(\model_\destdom)
\end{eqnarray*}


\subsection{Example Two: The TKF Structure Tree (TKFST)}

Consider now the RNA evolutionary model derived in 
\cite{Holmes2004}.

That model has stem ($S$) and loop ($L$) sequences, each of which is a TKF91 sequence with its own rates.
Stems are sequences of base pairs, loops of individual bases.

The model was developed as a proof of concept, but suffers similar deficiencies to TKF91 concerning the absence of an affine gap penalty,
as well as a lack of basepair stacking effects or other empirically observed features of biological RNA structures.

\newcommand\ruletable[3]{
\begin{tabular}{rcc|rrr}
\multicolumn{6}{l}{\em #1} \\
\multicolumn{6}{l}{\em #2}\\
#3
\end{tabular}
}
\newcommand\tkfstruletable[3]{\ruletable{TKF Structure Tree #1}{#2}{#3}}

\newcommand\lhsrhs[5]{ #1 & $\to$ & #2 & #3 & #4 & #5 \\ }
\newcommand\rhs[4]{       &   $|$ & #1 & #2 & #3 & #4 \\ }
\newcommand\blank{ & & & & & \\ }


\newcommand\singletgrammar[3]{
% #1 = tree node identifier
% #2 = 5' terminal (left)
% #3 = 3' terminal (right)
\tkfstruletable{singlet rules ($#1$)}
{Sequence $#1$ terminals: $\{ #2, #3 \}$.}
{
\lhsrhs{ lhs }{ rhs }{ $P(#1)$ }{  }{  }
\hline

\blank
\lhsrhs{ $L_{#1}$ }{ ${#2}\ L_{#1}$ }{ $\kappa_l \pi_l(#2)$ }{  }{  }
\rhs{ $S_{#1}\ L_{#1}$ }{ $\kappa_l \pi_l(S)$ }{  }{  }

\rhs{ $\epsilon$ }{ $1-\kappa_l$ }{  }{  }

\blank
\lhsrhs{ $S_{#1}$ }{ ${#2}\ S_{#1}\ {#3}$ }{ $\kappa_s \pi_s(#2#3)$ }{  }{  }
\rhs{ $L_{#1}$ }{ $1-\kappa_s$ }{  }{  }
}
}  % end \singletgrammar



\newcommand\pairgrammar[7]{
% #1 = ancestor
% #2 = descendant
% #3 = ancestral 5' terminal (left)
% #4 = ancestral 3' terminal (right)
% #5 = descendant 5' terminal (left)
% #6 = descendant 3' terminal (right)
% #7 = time suffix (empty string "", apostrophe "'", branch superscript "^{...}")
\tkfstruletable{pair rules ($#1 \stackrel{t#7}{\to} #2$)}
{Sequence $#1$ terminals: $\{ #3, #4 \}$. \quad  Sequence $#2$ terminals: $\{ #5, #6 \}$.}
{
\lhsrhs{ lhs }{ rhs }{ $P(#1)$ }{ $P(#2|#1)$ }{  }
\hline

\blank
\lhsrhs{ $L_{#1#2}$ }{ ${#3}\ {#5}\ L_{#1#2}$ }{ $\kappa_l \pi_l(#3)$ }{ $(1-\beta#7_l) \alpha#7_l M#7_l(#3,#5)$ }{  }
\rhs{ ${#5}\ L_{#1#2}$ }{ $1$ }{ $\beta#7_l \pi_l(#5)$ }{  }
\rhs{ ${#3}\ L_{#1\gap #2}$ }{ $\kappa_l \pi_l(#3)$ }{ $(1-\beta#7_l) (1-\alpha#7_l)$ }{  }

\blank
\rhs{ $S_{#1#2}\ L_{#1#2}$ }{ $\kappa_l \pi_l(S)$ }{ $(1-\beta#7_l) \alpha#7_l$ }{  }
\rhs{ $S_{#2}\ L_{#1#2}$ }{ $1$ }{ $\beta#7_l \pi_l(S)$ }{  }
\rhs{ $S_{#1}\ L_{#1\gap #2}$ }{ $\kappa_l \pi_l(S)$ }{ $(1-\beta#7_l) (1-\alpha#7_l)$ }{  }

\rhs{ $\epsilon$ }{ $1-\kappa_l$ }{ $1-\beta#7_l$ }{  }

\blank
\lhsrhs{ $S_{#1#2}$ }{ ${#3}\ {#5}\ S_{#1#2}\ {#6}\ {#4}$ }{ $\kappa_s \pi_s(#3#4)$ }{ $(1-\beta#7_s)  \alpha#7_s M#7_s(#3#4,#5#6)$ }{  }
\rhs{ ${#5}\ S_{#1#2}\ {#6}$ }{ $1$ }{ $\beta#7_s \pi_s(#5#6)$ }{  }
\rhs{ ${#3}\ S_{#1\gap #2}\ {#4}$ }{ $\kappa_s \pi_s(#3#4)$ }{ $(1-\beta#7_s) (1-\alpha#7_s)$ }{  }
\rhs{ $L_{#1#2}$ }{ $1-\kappa_s$ }{ $1-\beta#7_s$ }{  }

\blank
\lhsrhs{ $L_{#1\gap #2}$ }{ ${#3}\ {#5}\ L_{#1#2}$ }{ $\kappa_l \pi_l(#3)$ }{ $(1-\gamma#7_l) \alpha#7_l M#7_l(#3,#5)$ }{  }
\rhs{ ${#5}\ L_{#1#2}$ }{ $1$ }{ $\gamma#7_l \pi_l(#5)$ }{  }
\rhs{ ${#3}\ L_{#1\gap #2}$ }{ $\kappa_l \pi_l(#3)$ }{ $(1-\gamma#7_l) (1-\alpha#7_l)$ }{  }

\blank
\rhs{ $S_{#1#2}\ L_{#1#2}$ }{ $\kappa_l \pi_l(S)$ }{ $(1-\gamma#7_l) \alpha#7_l$ }{  }
\rhs{ $S_{#2}\ L_{#1#2}$ }{ $1$ }{ $\gamma#7_l \pi_l(S)$ }{  }
\rhs{ $S_{#1}\ L_{#1\gap #2}$ }{ $\kappa_l \pi_l(S)$ }{ $(1-\gamma#7_l) (1-\alpha#7_l)$ }{  }

\rhs{ $\epsilon$ }{ $1-\kappa_l$ }{ $1-\gamma#7_l$ }{  }

\blank
\lhsrhs{ $S_{#1\gap #2}$ }{ ${#3}\ {#5}\ S_{#1#2}\ {#6}\ {#4}$ }{ $\kappa_s \pi_s(#3#4)$ }{ $(1-\gamma#7_s) \alpha#7_s M#7_s(#3#4,#5#6)$ }{  }
\rhs{ ${#5}\ S_{#1#2}\ {#6}$ }{ $1$ }{ $\gamma#7_s \pi_s(#5#6)$ }{  }
\rhs{ ${#3}\ S_{#1\gap #2}\ {#4}$ }{ $\kappa_s \pi_s(#3#4)$ }{ $(1-\gamma#7_s) (1-\alpha#7_s)$ }{  }
\rhs{ $L_{#1#2}$ }{ $1-\kappa_s$ }{ $1-\gamma#7_s$ }{  }
}
}  % end \pairgrammar




\newcommand\abc{$a \stackrel{t}{\to} b \stackrel{t'}{\to} c$}
\newcommand\ab{$a \stackrel{t}{\to} b$}
\newcommand\bc{$b \stackrel{t'}{\to} c$}
\newcommand\ac{$a \stackrel{t+t'}{\longrightarrow} c$}

\newcommand\agapbc{a\gap bc}
\newcommand\abgapc{ab\gap c}

\newcommand\abcstates{``$abc$'' states (post emission at $c$)}
\newcommand\abgapcstates{``$\abgapc$'' states (post deletion on \bc\ branch)}
\newcommand\agapbcstates{``$\agapbc$'' states (post deletion on \ab\ branch)}


\begin{table}
\begin{center}
\singletgrammar{a}{u}{v}
\caption{TKF Structure Tree. Singlet rule-set for $a$.}
\end{center}
\end{table}
% b
\begin{table}
\begin{center}
\singletgrammar{b}{w}{x}
\caption{TKF Structure Tree. Singlet rule-set for $b$.}
\end{center}
\end{table}

\begin{table}
\begin{center}
\pairgrammar{a}{b}{u}{v}{w}{x}{}
\caption{TKF Structure Tree. Pair rule-set for \ab\ branch. Requires singlet rule-sets for $a$ and $b$.}
\end{center}
\end{table}



% Commands to build grammar tables
\newcommand\ntf[1]{\mbox{\tt #1}}  % flush nonterminals (no trailing space)
\newcommand\nt[1]{\ntf{#1}\ }  % nonterminals
\newcommand\ntm[1]{\ntf{#1}_M\ }  % match states
\newcommand\nti[1]{\ntf{#1}_I\ }  % insert states
\newcommand\ntd[1]{\ntf{#1}_D\ }  % delete states
\newcommand\ntx[1]{\ntf{#1}_X\ }  % X-only states
\newcommand\nty[1]{\ntf{#1}_Y\ }  % Y-only states
\newcommand\ntj[1]{\ntf{#1}_j\ }  % j is M, X, or Y
\newcommand\ntk[1]{\ntf{#1}_k\ }  % k is M or D

\newcommand\gramtab[1]{
\[
\begin{array}{rcl|ll}
\mbox{LHS}   & \to & \mbox{RHS} & \mbox{Transition} & \mbox{Emission} \\ \hline
#1
\end{array}
\]
}

\subsubsection{Parameters}

Let $\overline{xyz}$ denote reverse complement e.g. $\overline{AAG}=CTT$.

Parameters: insertion and deletion rates $\lambda_F < \mu_F$, fragment extension probability $r_F$,
substitution rate matrix $Q_F$, equilibrium probability vector $q_F$ (so $q_x Q_F = 0$)
for $F \in \{ R, N, G, S, L, C \}$.
Splice donor/acceptor site distribution $q_{D1}$,  $q_{D2}$, $q_{A1}$, $q_{A2}$.
Region-type probabilities $p_G + p_N + p_S + p_C = 1$.
Intron probability $p_I$.

The $Q_N$, $Q_S$, $Q_L$ and $Q_C$ models should be strand-invariant, so e.g. $Q_N(x_1,x_2) = Q_N(\overline{x_1},\overline{x_2})$.

\subsubsection{Functions}

For $F \in \{ R, N, G, S, L, C \}$:
\begin{eqnarray*}
\kappa_F & = & \left( 1 - \frac{\lambda_F}{\mu_F} \right) \left( 1 - r_F \right)
\end{eqnarray*}


\subsubsection{The Stationary Grammar}

\gramtab{
\nt{GENOME}  & \to & \nt{REGION} \nt{GENOME} & 1 - \kappa_R \\
             &   | & \epsilon & \kappa_R \\
\nt{REGION}  & \to & \nt{INTER} & p_N \\
             &   | & \nt{FWDCDS} & p_G / 2 \\
             &   | & \nt{REVCDS} & p_G / 2 \\
             &   | & \nt{STRUCT} & p_S \\
             &   | & \nt{CONS} & p_C \\
\nt{INTER}   & \to & x_1\ \nt{INTER} & 1 - \kappa_N & q_N(x_1) \\
             &   | & \epsilon & \kappa_N \\
\nt{FWDCDS}  & \to & \nt{FWDCOD} \nt{FWDCDS} & 1 - \kappa_G \\
             &   | & \epsilon & \kappa_G \\
\nt{FWDCOD}  & \to & x_1\ x_2\ x_3 & 1 - p_I & q_G(xyz) \\
             &   | & x_1\ x_2\ x_3\ \nt{FWDINT} & p_I / 3 & q_G(xyz) \\
             &   | & x_1\ x_2\ \nt{FWDINT} x_3 & p_I / 3 & q_G(xyz) \\
             &   | & x_1\ \nt{FWDINT} x_2\ x_3 & p_I / 3 & q_G(xyz) \\
\nt{FWDINT}  & \to & x_1\ x_2\ \nt{GENOME} x_3\ x_4 & 1 & q_{D1}(x_1) q_{D2}(x_2) q_{A1}(x_3) q_{A2}(x_4) \\
\nt{REVCDS}  & \to & \nt{REVCDS} \nt{REVCOD} & 1 - \kappa_G \\
             &   | & \epsilon & \kappa_G \\
\nt{REVCOD}  & \to & x_1\ x_2\ x_3 & 1 - p_I & q_G(\overline{xyz}) \\
             &   | & x_1\ x_2\ x_3\ \nt{REVINT} & p_I / 3 & q_G(\overline{xyz}) \\
             &   | & x_1\ x_2\ \nt{REVINT} x_3 & p_I / 3 & q_G(\overline{xyz}) \\
             &   | & x_1\ \nt{REVINT} x_2\ x_3 & p_I / 3 & q_G(\overline{xyz}) \\
\nt{REVINT}  & \to & x_1\ x_2\ \nt{GENOME} x_3\ x_4 & 1 & q_{D1}(\overline{x_4}) q_{D2}(\overline{x_3}) q_{A1}(\overline{x_2}) q_{A2}(\overline{x_1}) \\
\nt{STRUCT}  & \to & x_1\ \nt{STRUCT} x_2 & 1 - \kappa_S & q_S(xy) \\
             &   | & \nt{LOOP} & \kappa_S \\
\nt{LOOP}    & \to & x_1\ \nt{LOOP} & 1 - \kappa_L & q_L(x_1) \\
             &   | & \epsilon & \kappa_L \\
\nt{CONS}    & \to & x_1\ \nt{CONS} & 1 - \kappa_C & q_C(x_1) \\
             &   | & \epsilon & \kappa_C \\
}


\subsubsection{The Joint Finite-Time Grammar}

The general rules for forming the joint grammar from the conditional grammar are as follows.
For every nonterminal of the following form (here $\ntf{L}$, $\ntf{R}$, and/or $\ntf{E}$ are allowed to be $\epsilon$)
\gramtab{
\nt{F} & \to & \nt{L} \nt{F} \nt{R} & 1 - \kappa_F & q_F(LR) \\
       &   | & \nt{E} & \kappa_F \\
}
...that is, for $\ntf{GENOME}$, $\ntf{INTER}$, $\ntf{FWDCDS}$, $\ntf{REVCDS}$, $\ntf{STRUCT}$, $\ntf{LOOP}$, and $\ntf{CONS}$
(with $F \in \{ R, N, G, S, L, C \}$),
replace these rules with
\gramtab{
\ntm{F} & \to & \ntm{L} \ntm{F} \ntm{R} & (1 - \beta_{M,F}) (1 - \kappa_F) \alpha_F & q_F(LR_x) \exp(Q_F t)(LR_x,LR_y) \\
        &   | & \nty{L} \ntm{F} \nty{R} & \beta_{M,F} & q_F(LR_y) \\
        &   | & \ntx{L} \ntd{F} \ntx{R} & (1 - \beta_{M,F}) (1 - \kappa_F) (1 - \alpha_F) & q_F(LR_x) \\
        &   | & \ntm{E} & (1 - \beta_{M,F}) \kappa_F \\
\ntd{F} & \to & \ntm{L} \ntm{F} \ntm{R} & (1 - \beta_{D,F}) (1 - \kappa_F) \alpha_F & q_F(LR_x) \exp(Q_F t)(LR_x,LR_y) \\
        &   | & \nty{L} \ntm{F} \nty{R} & \beta_{D,F} & q_F(LR_y)  \\
        &   | & \ntx{L} \ntd{F} \ntx{R} & (1 - \beta_{D,F}) (1 - \kappa_F) (1 - \alpha_F) & q_F(LR_x) \\
        &   | & \ntm{E} & (1 - \beta_{D,F}) \kappa_F \\
}
...that is, two versions $\ntm{F}$ and $\ntd{F}$, with different $\beta$'s for each type.

For every other nonterminal $\ntf{N}$, there need to be three versions $\ntm{N}$, $\ntx{N}$ and $\nty{N}$,
with outgoing rules for each type going to other nonterminals of the same type.

\paragraph{Top level.}

\gramtab{
\ntk{GENOME}   & \to & \ntm{REGION} \ntm{GENOME} & (1 - \beta_{k,R}) (1 - \kappa_R) \alpha_R \\
               &   | & \nty{REGION} \ntm{GENOME} & \beta_{k,R}  \\
               &   | & \ntx{REGION} \ntd{GENOME} & (1 - \beta_{k,R}) (1 - \kappa_R) (1 - \alpha_R)  \\
               &   | & \epsilon & (1 - \beta_{k,R}) \kappa_R \\
\ntj{REGION}   & \to & \ntj{INTER} & p_N \\
               &   | & \ntj{FWDCDS} & p_G / 2 \\
               &   | & \ntj{REVCDS} & p_G / 2 \\
               &   | & \ntj{STRUCT} & p_S \\
               &   | & \ntj{CONS} & p_C \\
\nt{INTER}   & \to & x_1\ \nt{INTER} & 1 - \kappa_N & q_N(x_1) \\
             &   | & \epsilon & \kappa_N \\
}

\paragraph{Coding sequences.}

\gramtab{
\nt{FWDCDS}  & \to & \nt{FWDCOD} \nt{FWDCDS} & 1 - \kappa_G \\
             &   | & \epsilon & \kappa_G \\
\nt{FWDCOD}  & \to & x_1\ x_2\ x_3 & 1 - p_I & q_G(xyz) \\
             &   | & x_1\ x_2\ x_3\ \nt{FWDINT} & p_I / 3 & q_G(xyz) \\
             &   | & x_1\ x_2\ \nt{FWDINT} x_3 & p_I / 3 & q_G(xyz) \\
             &   | & x_1\ \nt{FWDINT} x_2\ x_3 & p_I / 3 & q_G(xyz) \\
\nt{FWDINT}  & \to & x_1\ x_2\ \nt{GENOME} x_3\ x_4 & 1 & q_{D1}(x_1) q_{D2}(x_2) q_{A1}(x_3) q_{A2}(x_4) \\
}

\gramtab{
\nt{REVCDS}  & \to & \nt{REVCDS} \nt{REVCOD} & 1 - \kappa_G \\
             &   | & \epsilon & \kappa_G \\
\nt{REVCOD}  & \to & x_1\ x_2\ x_3 & 1 - p_I & q_G(\overline{xyz}) \\
             &   | & x_1\ x_2\ x_3\ \nt{REVINT} & p_I / 3 & q_G(\overline{xyz}) \\
             &   | & x_1\ x_2\ \nt{REVINT} x_3 & p_I / 3 & q_G(\overline{xyz}) \\
             &   | & x_1\ \nt{REVINT} x_2\ x_3 & p_I / 3 & q_G(\overline{xyz}) \\
\nt{REVINT}  & \to & x_1\ x_2\ \nt{GENOME} x_3\ x_4 & 1 & q_{D1}(\overline{x_4}) q_{D2}(\overline{x_3}) q_{A1}(\overline{x_2}) q_{A2}(\overline{x_1}) \\
}

\paragraph{RNA structures.}

\gramtab{
\nt{STRUCT}  & \to & x_1\ \nt{STRUCT} x_2 & 1 - \kappa_S & q_S(xy) \\
             &   | & \nt{LOOP} & \kappa_S \\
\nt{LOOP}    & \to & x_1\ \nt{LOOP} & 1 - \kappa_L & q_L(x_1) \\
             &   | & \epsilon & \kappa_L \\
\nt{CONS}    & \to & x_1\ \nt{CONS} & 1 - \kappa_C & q_C(x_1) \\
             &   | & \epsilon & \kappa_C \\
}


\subsection{Example Three: The TKF Basepair Stack (TKFStack)}
\label{sec:rna-evol}

The simple TKF Structure Tree models RNA secondary structure with
alternating stems (basepair sequences) and loops (single-nucleotide
sequences). While a useful proof of concept, it lacks basepair stacking,
multiloop junctions, bulges, and internal loops---features critical for
realistic RNA structure modeling.

We now define an enhanced stem-loop grammar that incorporates these
features while remaining within the TKF evolutionary framework. We then
show how evolution elaboration, profile SCFG construction, and the
triplet model for progressive reconstruction all apply to this grammar.

\subsubsection{Parameters}

The grammar has the following parameters:
\begin{itemize}
\item Stem link TKF91 rates $\insrate_S < \delrate_S$,
  giving $\kappa_S = \insrate_S / \delrate_S$.
\item Loop link TKF91 rates $\insrate_L < \delrate_L$,
  giving $\kappa_L = \insrate_L / \delrate_L$.
\item Bulge extension probability $\ext_B$.
\item Stacked-pair fragment extension probability $\ext_K$.
\item Stem link type probabilities
  $p_{\mathrm{bp}} + p_{\mathrm{st}} + p_{\mathrm{bu}} = 1$.
\item Loop link type probabilities
  $p_{\mathrm{lf}} + p_{\mathrm{rf}} + p_{\mathrm{sl}} = 1$.
\item Nucleotide equilibrium $\eqm(c)$ over $\alphabet = \{A,C,G,U\}$,
  with rate matrix $\exch$.
\item Single basepair equilibrium $\eqm_{\mathrm{bp}}(c_L, c_R)$
  over $|\alphabet|^2 = 16$ states, rate matrix $\exch_{\mathrm{bp}}$.
\item Closing basepair equilibrium $\eqm_{\mathrm{cl}}(c_L, c_R)$
  over $16$ states, rate matrix $\exch_{\mathrm{cl}}$.
\item Stacked-pair equilibrium
  $\eqm_K(c_L^1, c_L^2, c_R^2, c_R^1)$
  over the $6^2 = 36$ canonical stacked pairs, rate matrix $\exch_K$.
  A stacked pair consists of two consecutive canonical basepairs
  $(c_L^1, c_R^1)$ and $(c_L^2, c_R^2)$; the state space is
  restricted to the six Watson-Crick and wobble pairs
  (AU, CG, GC, UA, GU, UG) for each position,
  giving 36 states rather than the full $16^2 = 256$.
\end{itemize}

\subsubsection{Nonterminals and Productions}

A \emph{stem-loop} consists of a stem (nested basepairs with
decorations), a closing basepair, and a loop (with possible
multiloop branches).
The start symbol is $\mathsf{SL}$ (stem-loop).

\medskip
\noindent\textbf{Stem-loop:}
\begin{align}
\mathsf{SL} &\to \mathsf{STEM}
  && \text{weight } 1
  \label{eq:sl}
\end{align}

\noindent\textbf{Stem (TKF91 sequence of stem-links, from outer to inner):}
\begin{align}
\mathsf{STEM} &\to c_L\;\mathsf{STEM}\;c_R
  && \kappa_S\, p_{\mathrm{bp}}\, \eqm_{\mathrm{bp}}(c_L, c_R)
  && \text{[single basepair, LR]}
  \label{eq:stem-bp} \\
&\to c_L^1\, c_L^2\;\mathsf{STEM}\;c_R^2\, c_R^1
  && \kappa_S\, p_{\mathrm{st}}\, (1-\ext_K)\, \eqm_K(\cdot)
  && \text{[terminal stacked pair, LLRR]}
  \label{eq:stem-stack-term} \\
&\to c_L^1\, c_L^2\;\mathsf{STACK}\;c_R^2\, c_R^1
  && \kappa_S\, p_{\mathrm{st}}\, \ext_K\, \eqm_K(\cdot)
  && \text{[extended stacked pair, LLRR]}
  \label{eq:stem-stack-ext} \\
&\to \mathsf{LDECO}\;\mathsf{STEM}\;\mathsf{RDECO}
  && \kappa_S\, p_{\mathrm{bu}}
  && \text{[bulge]}
  \label{eq:stem-bulge} \\
&\to \mathsf{CLOSE}\;\mathsf{LOOP}
  && 1 - \kappa_S
  && \text{[end stem]}
  \label{eq:stem-end}
\end{align}
A bulge link~\eqref{eq:stem-bulge} represents an internal loop,
a single-sided bulge, or a multihelix junction branch between
basepairs.
All non-basepair content between consecutive basepairs is
consolidated into a single TKF link type whose internal structure
is governed by $\mathsf{LDECO}$ and $\mathsf{RDECO}$.

\medskip
\noindent\textbf{Bulge decorations (L-side and R-side content):}
\begin{align}
\mathsf{LDECO} &\to \mathsf{LFRAG}\;\mathsf{LDECO}
  && \text{[L-fragment, then more L-content]}
  \label{eq:ldeco-lfrag} \\
&\to \mathsf{SL}\;\mathsf{LDECO}
  && \text{[left branch: nested stem-loop]}
  \label{eq:ldeco-branch} \\
&\to \epsilon
  && \text{[end L-decorations]}
  \label{eq:ldeco-end} \\[6pt]
\mathsf{RDECO} &\to \mathsf{RDECO}\;\mathsf{RFRAG}
  && \text{[R-fragment, more R-content]}
  \label{eq:rdeco-rfrag} \\
&\to \mathsf{RDECO}\;\mathsf{SL}
  && \text{[right branch: nested stem-loop]}
  \label{eq:rdeco-branch} \\
&\to \epsilon
  && \text{[end R-decorations]}
  \label{eq:rdeco-end}
\end{align}
$\mathsf{LDECO}$ generates all content on the $5'$ side
(between the outer basepair's L-half and the continuation),
while $\mathsf{RDECO}$ generates all content on the $3'$ side
(between the continuation and the outer basepair's R-half).
An internal loop has both $\mathsf{LDECO} \neq \epsilon$ and
$\mathsf{RDECO} \neq \epsilon$;
a single-sided bulge has content on only one side;
a multihelix junction branch has $\mathsf{SL}$ nested within
$\mathsf{LDECO}$ or $\mathsf{RDECO}$.

\medskip
\noindent\textbf{Stacked-pair fragment continuation (LLRR emission):}
\begin{align}
\mathsf{STACK} &\to c_L^1\, c_L^2\;\mathsf{STEM}\;c_R^2\, c_R^1
  && (1-\ext_K)\, \eqm_K(\cdot)
  && \text{[terminal]}
  \label{eq:stack-term} \\
&\to c_L^1\, c_L^2\;\mathsf{STACK}\;c_R^2\, c_R^1
  && \ext_K\, \eqm_K(\cdot)
  && \text{[extend]}
  \label{eq:stack-ext}
\end{align}

\noindent\textbf{Closing basepair (LR emission):}
\begin{align}
\mathsf{CLOSE} &\to c_L\;\mathsf{CLOSE}'\;c_R
  && \eqm_{\mathrm{cl}}(c_L, c_R)
  \label{eq:close}
\end{align}
where $\mathsf{CLOSE}'$ is a unit nonterminal ($\mathsf{CLOSE}' \to \epsilon$, weight~$1$)
that serves as a placeholder for the inside of the closing basepair.

\noindent\textbf{Loop (TKF91 sequence of loop-links):}
\begin{align}
\mathsf{LOOP} &\to \mathsf{LOOPLINK}\;\mathsf{LOOP}
  && \kappa_L
  && \text{[add loop link]}
  \label{eq:loop-link} \\
&\to \epsilon
  && 1 - \kappa_L
  && \text{[end loop]}
  \label{eq:loop-end}
\end{align}

\noindent\textbf{Loop link types:}
\begin{align}
\mathsf{LOOPLINK} &\to \mathsf{LFRAG}
  && p_{\mathrm{lf}}
  && \text{[L-fragment]}
  \label{eq:loop-lfrag} \\
&\to \mathsf{RFRAG}
  && p_{\mathrm{rf}}
  && \text{[R-fragment]}
  \label{eq:loop-rfrag} \\
&\to \mathsf{SL}
  && p_{\mathrm{sl}}
  && \text{[nested stem-loop (multiloop)]}
  \label{eq:loop-sl}
\end{align}

\noindent\textbf{Unpaired nucleotide fragments in loops
(geometric length $\geq 1$):}
\begin{align}
\mathsf{LFRAG} &\to c\;\mathsf{LFRAG}
  && \ext_L\, \eqm(c)
  && \text{[extend, L-emission]}
  \label{eq:lfrag-ext} \\
&\to c
  && (1-\ext_L)\, \eqm(c)
  && \text{[terminal, L-emission]}
  \label{eq:lfrag-term} \\[6pt]
\mathsf{RFRAG} &\to \mathsf{RFRAG}\;c
  && \ext_R\, \eqm(c)
  && \text{[extend, R-emission]}
  \label{eq:rfrag-ext} \\
&\to c
  && (1-\ext_R)\, \eqm(c)
  && \text{[terminal, R-emission]}
  \label{eq:rfrag-term}
\end{align}

\subsubsection{Emission Types and Span Tracking}
\label{sec:emission-types}

The grammar has four emission patterns, each determining how
terminals consume positions from the span $[i,j]$ of the
input sequence:

\begin{enumerate}
\item \textbf{L-emission} ($c$):
  consumes from the left end, advancing $i \to i+1$.
  Used by $\mathsf{LFRAG}$, $\mathsf{LDECO}$ content.

\item \textbf{R-emission} ($c$):
  consumes from the right end, retreating $j \to j-1$.
  Used by $\mathsf{RFRAG}$, $\mathsf{RDECO}$ content.

\item \textbf{LR-emission} ($c_L, c_R$):
  consumes from both ends simultaneously, $i \to i+1$, $j \to j-1$.
  Used by $\mathsf{STEM}$ (single basepair), $\mathsf{CLOSE}$.

\item \textbf{LLRR-emission} ($c_L^1, c_L^2, c_R^2, c_R^1$):
  consumes two from each end, $i \to i+2$, $j \to j-2$.
  Used by $\mathsf{STEM}$ (stacked pair), $\mathsf{STACK}$.
  Each LLRR unit represents a \emph{stacked dinucleotide pair}:
  two consecutive canonical basepairs $(c_L^1, c_R^1)$ and $(c_L^2, c_R^2)$
  treated as a single evolutionary unit with
  nearest-neighbor stacking energy.
  The state space is restricted to the $6 \times 6 = 36$
  combinations of canonical (Watson-Crick and wobble) pairs.
\end{enumerate}

\begin{remark}[L/R distinction in terminals vs.\ productions]
\label{rem:lr-distinction}
The terminal alphabet is $\alphabet = \{A,C,G,U\}$ without
L/R copies: both L-emission and R-emission produce characters
from the same alphabet, differing only in which end of the
span $[i,j]$ is consumed.
At leaf nodes, the observed sequence carries no L/R annotation.

However, the L/R distinction is essential at the \emph{production}
level and therefore in the \emph{profile SCFG}.
Each nonterminal instance in a profile has a fixed emission
direction (L, R, LR, or LLRR), determined by the parse tree
from which it was extracted---analogous to the MATL, MATR,
and MATP node types in Infernal covariance models~\cite{Nawrocki2009}.
The branch PTT must preserve emission-direction compatibility
when mapping parent productions to child productions.
Consequently, a profile position that is LR-emitting (basepaired)
at one node of the phylogeny may correspond to separate
L-emitting and R-emitting positions (unpaired) at another node,
reflecting structural change along the evolutionary lineage.
\end{remark}

\begin{remark}[Structural interpretation]
Consider a stem with links (from outer to inner):
basepair $(c_L^1, c_R^1)$, bulge (left-side $b_1 b_2$,
right-side branch $\mathsf{SL}'$), basepair $(c_L^2, c_R^2)$.
The yield is:
\[
c_L^1\; b_1\; b_2\; c_L^2\; [\text{loop}]\;
c_R^2\; [\text{yield}(\mathsf{SL}')]\; c_R^1
\]
The bulge nucleotides $b_1, b_2$ appear on the $5'$ side
(via $\mathsf{LDECO}$), while the branch structure appears
on the $3'$ side (via $\mathsf{RDECO}$).
This represents an internal loop with unpaired nucleotides
on the $5'$ strand and a branching sub-structure on the $3'$
strand---a configuration that the consolidated bulge link
type~\eqref{eq:stem-bulge} captures as a single TKF event.
\end{remark}

\subsubsection{The Pair Grammar via Evolution Elaboration}
\label{sec:pair-grammar}

Applying the evolution elaboration rules
(Appendix~\ref{sec:appendix-grammar}) to the singlet grammar
produces the pair grammar for an ancestor--descendant pair
separated by evolutionary time~$\evoltime$.
Each nonterminal participating in a TKF91 link sequence
($\mathsf{STEM}$, $\mathsf{LOOP}$) gains $\mat/\ins/\del$
versions with the standard TKF91 transition weights
$(\alpha, \beta, \gamma)$ derived from $(\insrate, \delrate, \evoltime)$.

\paragraph{Stem Link Elaboration.}

The stem is a TKF91 sequence of links with rates
$(\insrate_S, \delrate_S)$.
Each $\mathsf{STEM}$ nonterminal becomes
$\mathsf{STEM}_\mat$ (post-match/insert) and
$\mathsf{STEM}_\del$ (post-delete):

\begin{align}
\mathsf{STEM}_\mat &\to c_L^x\, c_L^y\;\mathsf{STEM}_\mat\;c_R^y\, c_R^x
  && (1-\beta_S) \kappa_S\, p_{\mathrm{bp}}\, \alpha_S\,
     \eqm_{\mathrm{bp}}(c_L^x, c_R^x)\,
     P_{\mathrm{bp}}(c_L^y, c_R^y | c_L^x, c_R^x)
  \label{eq:pair-bp-match} \\
&\to c_L^y\;\mathsf{STEM}_\mat\;c_R^y
  && \beta_S\, \kappa_S\, p_{\mathrm{bp}}\,
     \eqm_{\mathrm{bp}}(c_L^y, c_R^y)
  && \text{[BP insert]}
  \label{eq:pair-bp-ins} \\
&\to c_L^x\;\mathsf{STEM}_\del\;c_R^x
  && (1-\beta_S) \kappa_S\, p_{\mathrm{bp}}\, (1-\alpha_S)\,
     \eqm_{\mathrm{bp}}(c_L^x, c_R^x)
  && \text{[BP delete]}
  \label{eq:pair-bp-del}
\end{align}
with analogous rules for $\mathsf{STEM}_\del$ using
$\gamma_S$ in place of $\beta_S$.

For matched single basepairs~\eqref{eq:pair-bp-match}, the nesting
is $c_L^x\, c_L^y\;\cdots\;c_R^y\, c_R^x$: ancestor terminals
on the outside, descendant terminals on the inside,
preserving palindromic structure.
For inserted basepairs~\eqref{eq:pair-bp-ins}, only descendant
terminals appear (LR emission).
For deleted basepairs~\eqref{eq:pair-bp-del}, only ancestor
terminals appear (LR emission).

\paragraph{Stacked-Pair Elaboration.}

Matched stacked pairs have LLRR emission on both ancestor
and descendant, yielding an $L^4R^4$ nesting pattern
in the pair grammar:
\begin{align}
\mathsf{STEM}_\mat &\to
  c_{L1}^x\, c_{L2}^x\, c_{L1}^y\, c_{L2}^y\;
  \mathsf{STEM}_\mat\;
  c_{R2}^y\, c_{R1}^y\, c_{R2}^x\, c_{R1}^x
  \notag \\
  &\qquad\qquad
  (1-\beta_S) \kappa_S\, p_{\mathrm{st}}\, (1-\ext_K)\, \alpha_S\,
  \eqm_K(\cdot)\, P_K(\cdot|\cdot)
  \label{eq:pair-stack-match}
\end{align}
with inserted stacked pairs emitting only descendant LLRR,
and deleted stacked pairs emitting only ancestor LLRR.

\paragraph{Bulge Elaboration.}

A bulge link~\eqref{eq:stem-bulge} gains $\mat/\ins/\del$
versions like any other stem link.
A matched bulge has the form
$\mathsf{LDECO}_\mat\;\mathsf{STEM}_\mat\;\mathsf{RDECO}_\mat$,
with each side elaborated independently.

Within $\mathsf{LDECO}$, each sub-element
($\mathsf{LFRAG}$, nested $\mathsf{SL}$) gains $\mat/\ins/\del$
versions with L-emission direction preserved.
For example, a matched left-branch within a bulge produces:
\begin{align}
\mathsf{LDECO}_\mat &\to \mathsf{SL}_\mat\;\mathsf{LDECO}_\mat
  && \alpha_{\mathrm{deco}}\, P(\text{branch})
  && \text{[matched left branch]}
  \label{eq:pair-ldeco-match} \\
&\to \mathsf{SL}_\ins\;\mathsf{LDECO}_\mat
  && \beta_{\mathrm{deco}}
  && \text{[inserted left branch]}
  \label{eq:pair-ldeco-ins} \\
&\to \mathsf{SL}_\del\;\mathsf{LDECO}_\del
  && (1-\alpha_{\mathrm{deco}})\, P(\text{branch})
  && \text{[deleted left branch]}
  \label{eq:pair-ldeco-del}
\end{align}
Here $\mathsf{SL}_\mat$ generates aligned ancestor--descendant
sub-structures, $\mathsf{SL}_\ins$ generates descendant-only
sub-structures, and $\mathsf{SL}_\del$ generates ancestor-only
sub-structures.

$\mathsf{RDECO}$ is elaborated symmetrically, with R-emission
direction preserved for $\mathsf{RFRAG}$ elements and right-side
branches.

\paragraph{Loop Elaboration.}

The loop is a TKF91 sequence of loop-links with rates
$(\insrate_L, \delrate_L)$.
Elaboration proceeds identically to the stem: each
$\mathsf{LOOP}$ becomes $\mathsf{LOOP}_\mat / \mathsf{LOOP}_\del$,
and each loop-link type ($\mathsf{LFRAG}$, $\mathsf{RFRAG}$,
nested $\mathsf{SL}$) gains $\mat/\ins/\del$ versions.

Nested stem-loops within the loop (multiloop junctions) are
handled recursively: $\mathsf{SL}_\mat$ aligns both ancestor
and descendant sub-structures, allowing multiloop branches to
be independently inserted, deleted, or matched.


\subsection{Example Four: The TKF Genome}


\subsubsection{Parameters}

Let $\overline{xyz}$ denote reverse complement e.g. $\overline{AAG}=CTT$.

Parameters: insertion and deletion rates $\lambda_F < \mu_F$, fragment extension probability $r_F$,
substitution rate matrix $Q_F$, equilibrium probability vector $q_F$ (so $q_x Q_F = 0$)
for $F \in \{ R, N, G, S, L, C \}$.
Splice donor/acceptor site distribution $q_{D1}$,  $q_{D2}$, $q_{A1}$, $q_{A2}$.
Region-type probabilities $p_G + p_N + p_S + p_C = 1$.
Intron probability $p_I$.

The $Q_N$, $Q_S$, $Q_L$ and $Q_C$ models should be strand-invariant, so e.g. $Q_N(x_1,x_2) = Q_N(\overline{x_1},\overline{x_2})$.

\subsubsection{Functions}

For $F \in \{ R, N, G, S, L, C \}$:
\begin{eqnarray*}
\kappa_F & = & \left( 1 - \frac{\lambda_F}{\mu_F} \right) \left( 1 - r_F \right)
\end{eqnarray*}


\subsubsection{The Stationary Grammar}

\gramtab{
\nt{GENOME}  & \to & \nt{REGION} \nt{GENOME} & 1 - \kappa_R \\
             &   | & \epsilon & \kappa_R \\
\nt{REGION}  & \to & \nt{INTER} & p_N \\
             &   | & \nt{FWDCDS} & p_G / 2 \\
             &   | & \nt{REVCDS} & p_G / 2 \\
             &   | & \nt{STRUCT} & p_S \\
             &   | & \nt{CONS} & p_C \\
\nt{INTER}   & \to & x_1\ \nt{INTER} & 1 - \kappa_N & q_N(x_1) \\
             &   | & \epsilon & \kappa_N \\
\nt{FWDCDS}  & \to & \nt{FWDCOD} \nt{FWDCDS} & 1 - \kappa_G \\
             &   | & \epsilon & \kappa_G \\
\nt{FWDCOD}  & \to & x_1\ x_2\ x_3 & 1 - p_I & q_G(xyz) \\
             &   | & x_1\ x_2\ x_3\ \nt{FWDINT} & p_I / 3 & q_G(xyz) \\
             &   | & x_1\ x_2\ \nt{FWDINT} x_3 & p_I / 3 & q_G(xyz) \\
             &   | & x_1\ \nt{FWDINT} x_2\ x_3 & p_I / 3 & q_G(xyz) \\
\nt{FWDINT}  & \to & x_1\ x_2\ \nt{GENOME} x_3\ x_4 & 1 & q_{D1}(x_1) q_{D2}(x_2) q_{A1}(x_3) q_{A2}(x_4) \\
\nt{REVCDS}  & \to & \nt{REVCDS} \nt{REVCOD} & 1 - \kappa_G \\
             &   | & \epsilon & \kappa_G \\
\nt{REVCOD}  & \to & x_1\ x_2\ x_3 & 1 - p_I & q_G(\overline{xyz}) \\
             &   | & x_1\ x_2\ x_3\ \nt{REVINT} & p_I / 3 & q_G(\overline{xyz}) \\
             &   | & x_1\ x_2\ \nt{REVINT} x_3 & p_I / 3 & q_G(\overline{xyz}) \\
             &   | & x_1\ \nt{REVINT} x_2\ x_3 & p_I / 3 & q_G(\overline{xyz}) \\
\nt{REVINT}  & \to & x_1\ x_2\ \nt{GENOME} x_3\ x_4 & 1 & q_{D1}(\overline{x_4}) q_{D2}(\overline{x_3}) q_{A1}(\overline{x_2}) q_{A2}(\overline{x_1}) \\
\nt{STRUCT}  & \to & x_1\ \nt{STRUCT} x_2 & 1 - \kappa_S & q_S(xy) \\
             &   | & \nt{LOOP} & \kappa_S \\
\nt{LOOP}    & \to & x_1\ \nt{LOOP} & 1 - \kappa_L & q_L(x_1) \\
             &   | & \epsilon & \kappa_L \\
\nt{CONS}    & \to & x_1\ \nt{CONS} & 1 - \kappa_C & q_C(x_1) \\
             &   | & \epsilon & \kappa_C \\
}


\subsubsection{The Joint Finite-Time Grammar}

\paragraph{How to form the joint grammar}
The general rules for forming the joint grammar from the conditional grammar are as follows.
For every nonterminal of the following form (here $\ntf{L}$, $\ntf{R}$, and/or $\ntf{E}$ are allowed to be $\epsilon$)
\gramtab{
\nt{F} & \to & \nt{L} \nt{F} \nt{R} & 1 - \kappa_F & q_F(LR) \\
       &   | & \nt{E} & \kappa_F \\
}
...that is, for $\ntf{GENOME}$, $\ntf{INTER}$, $\ntf{FWDCDS}$, $\ntf{REVCDS}$, $\ntf{STRUCT}$, $\ntf{LOOP}$, and $\ntf{CONS}$
(with $F \in \{ R, N, G, S, L, C \}$),
replace these rules with
\gramtab{
\ntm{F} & \to & \ntm{L} \ntm{F} \ntm{R} & (1 - \beta_{M,F}) (1 - \kappa_F) \alpha_F & q_F(LR_x) \exp(Q_F t)(LR_x,LR_y) \\
        &   | & \nty{L} \ntm{F} \nty{R} & \beta_{M,F} & q_F(LR_y) \\
        &   | & \ntx{L} \ntd{F} \ntx{R} & (1 - \beta_{M,F}) (1 - \kappa_F) (1 - \alpha_F) & q_F(LR_x) \\
        &   | & \ntm{E} & (1 - \beta_{M,F}) \kappa_F \\
\ntd{F} & \to & \ntm{L} \ntm{F} \ntm{R} & (1 - \beta_{D,F}) (1 - \kappa_F) \alpha_F & q_F(LR_x) \exp(Q_F t)(LR_x,LR_y) \\
        &   | & \nty{L} \ntm{F} \nty{R} & \beta_{D,F} & q_F(LR_y)  \\
        &   | & \ntx{L} \ntd{F} \ntx{R} & (1 - \beta_{D,F}) (1 - \kappa_F) (1 - \alpha_F) & q_F(LR_x) \\
        &   | & \ntm{E} & (1 - \beta_{D,F}) \kappa_F \\
}
...that is, two versions $\ntm{F}$ and $\ntd{F}$, with different $\beta$'s for each type.

For every other nonterminal $\ntf{N}$, there need to be three versions $\ntm{N}$, $\ntx{N}$ and $\nty{N}$,
with outgoing rules for each type going to other nonterminals of the same type.

\paragraph{Top level}

\gramtab{
\ntk{GENOME}   & \to & \ntm{REGION} \ntm{GENOME} & (1 - \beta_{k,R}) (1 - \kappa_R) \alpha_R \\
               &   | & \nty{REGION} \ntm{GENOME} & \beta_{k,R}  \\
               &   | & \ntx{REGION} \ntd{GENOME} & (1 - \beta_{k,R}) (1 - \kappa_R) (1 - \alpha_R)  \\
               &   | & \epsilon & (1 - \beta_{k,R}) \kappa_R \\
\ntj{REGION}   & \to & \ntj{INTER} & p_N \\
               &   | & \ntj{FWDCDS} & p_G / 2 \\
               &   | & \ntj{REVCDS} & p_G / 2 \\
               &   | & \ntj{STRUCT} & p_S \\
               &   | & \ntj{CONS} & p_C \\
\nt{INTER}   & \to & x_1\ \nt{INTER} & 1 - \kappa_N & q_N(x_1) \\
             &   | & \epsilon & \kappa_N \\
}

\paragraph{Coding sequences}

\gramtab{
\nt{FWDCDS}  & \to & \nt{FWDCOD} \nt{FWDCDS} & 1 - \kappa_G \\
             &   | & \epsilon & \kappa_G \\
\nt{FWDCOD}  & \to & x_1\ x_2\ x_3 & 1 - p_I & q_G(xyz) \\
             &   | & x_1\ x_2\ x_3\ \nt{FWDINT} & p_I / 3 & q_G(xyz) \\
             &   | & x_1\ x_2\ \nt{FWDINT} x_3 & p_I / 3 & q_G(xyz) \\
             &   | & x_1\ \nt{FWDINT} x_2\ x_3 & p_I / 3 & q_G(xyz) \\
\nt{FWDINT}  & \to & x_1\ x_2\ \nt{GENOME} x_3\ x_4 & 1 & q_{D1}(x_1) q_{D2}(x_2) q_{A1}(x_3) q_{A2}(x_4) \\
}

\gramtab{
\nt{REVCDS}  & \to & \nt{REVCDS} \nt{REVCOD} & 1 - \kappa_G \\
             &   | & \epsilon & \kappa_G \\
\nt{REVCOD}  & \to & x_1\ x_2\ x_3 & 1 - p_I & q_G(\overline{xyz}) \\
             &   | & x_1\ x_2\ x_3\ \nt{REVINT} & p_I / 3 & q_G(\overline{xyz}) \\
             &   | & x_1\ x_2\ \nt{REVINT} x_3 & p_I / 3 & q_G(\overline{xyz}) \\
             &   | & x_1\ \nt{REVINT} x_2\ x_3 & p_I / 3 & q_G(\overline{xyz}) \\
\nt{REVINT}  & \to & x_1\ x_2\ \nt{GENOME} x_3\ x_4 & 1 & q_{D1}(\overline{x_4}) q_{D2}(\overline{x_3}) q_{A1}(\overline{x_2}) q_{A2}(\overline{x_1}) \\
}

\paragraph{RNA structures}

\gramtab{
\nt{STRUCT}  & \to & x_1\ \nt{STRUCT} x_2 & 1 - \kappa_S & q_S(xy) \\
             &   | & \nt{LOOP} & \kappa_S \\
\nt{LOOP}    & \to & x_1\ \nt{LOOP} & 1 - \kappa_L & q_L(x_1) \\
             &   | & \epsilon & \kappa_L \\
\nt{CONS}    & \to & x_1\ \nt{CONS} & 1 - \kappa_C & q_C(x_1) \\
             &   | & \epsilon & \kappa_C \\
}
